\documentclass[12pt,a4paper]{article}

%% Para acentos
\usepackage[latin1]{inputenc}

% Para fuentes postscript
%\usepackage{type1ec}
%\usepackage{lmodern}
\usepackage[T1]{fontenc}


%% Para mates
\usepackage{amsmath,amssymb,latexsym}

%% Para gr�ficos
\usepackage{graphicx}

%% Para gr�ficos realizados con LatexCad
%\usepackage{latexcad}


%% Para numeraci�n autom�tica
\newcounter{nroproblema}
\setcounter{nroproblema}{1}
\newcounter{nroapartado}
\setcounter{nroapartado}{1}
\newcommand{\n}{%
    \vspace{.1in}\noindent{\bf \arabic{nroproblema})\ }%
    \addtocounter{nroproblema}{1}%
    \setcounter{nroapartado}{1}%
    }%
\newcommand{\apartado}{\vspace*{0.2cm}\hspace*{0.2cm}%
   {\bf \alph{nroapartado})}
   \addtocounter{nroapartado}{1}%
   }


%% Abreviaturas comunes (para evitar usar misMacros.mcr)
\newcommand{\ds}{\displaystyle}
\newcommand{\df}{\mathop=\limits^{\sf def}}

% Para maquetaci�n
\setlength{\parindent}{0pt}
\usepackage[hmargin={1.5cm,1.5cm},vmargin={1.5cm,1.5cm}]{geometry}



\begin{document}



\begin{center}{\bf {\Large \'ALGEBRA I. HOJA 8}}\end{center}


%\begin{enumerate}

    \n Calcular  los determinantes
    siguientes:
    $$
a)\left|%
\begin{array}{ccc}
  0 & 2 & 4 \\
  1 & 1 & 3 \\
  3 & 3 & 3 \\
\end{array}%
\right|\ b)\left|%
\begin{array}{ccc}
  0 & -1 & 1 \\
  1 & 0 & 1 \\
  2 & -1 & -1 \\
\end{array}%
\right|\ c)\left|%
\begin{array}{ccc}
  -2 & 1 & 1 \\
  -3 & 2 & 1 \\
  1 & -1 & 1 \\
\end{array}%
\right|\ d)\left|%
\begin{array}{cccc}
  0 & -1 & -1 & 1 \\
  1 & 0 & -1 & 1 \\
  2 & -1 & -3 & -1 \\
  0 & 1 & -1 & 0 \\
\end{array}%
\right|
$$

Sol: 12, -4, -1, -8.
    
    \n Calcular los determinantes de las matrices dadas a
    continuaci�n:
    $$
a)\left|\frac{1}{7}\left(%
\begin{array}{ccc}
  3 & 2 & -6 \\
  6 & -3 & 2 \\
  2 & 6 & 3 \\
\end{array}%
\right) \right|
\hspace{1cm} b)\left|\frac{1}{9}\left(%
\begin{array}{ccc}
  -8 & -1 & -4 \\
  -4 & 4 & 7 \\
  1 & 8 & -4 \\
\end{array}%
\right) \right|
    $$
    
    Sol: a) -1,  b) 1.
    
    \n Demostrar que los
    siguientes determinantes son nulos:
    $$
a)\left|\begin{array}{ccc}
  7 & 1 & 3 \\
  3 & 2 & 1 \\
  10 & 3 & 4 \\
\end{array}\right|\hspace{5mm}
b)\left|\begin{array}{ccc}
  a & b & c \\
  -a & -b & -c \\
  d & e & f \\
\end{array}\right|\hspace{5mm}
c)\left|\begin{array}{ccc}
  1 & 2 & 3 \\
  2 & 3 & 4 \\
  3 & 4 & 5 \\
\end{array}\right|\hspace{5mm}
d)\left|\begin{array}{cccc}
  1^2 & 2^2 & 3^2 & 4^2 \\
  2^2 & 3^2 & 4^2 & 5^2 \\
  3^2 & 4^2 & 5^2 & 6^2 \\
  4^2 & 5^2 & 6^2 & 7^2 \\
\end{array}\right|
    $$
    
    \n
    a) Comprobar que la ecuaci�n de un plano que pasa por tres
    puntos $(a_1,a_2,a_3)$, $(b_1,b_2,b_3)$, $(c_1,c_2,c_3)$ del
    espacio es
    $$
\left|\begin{array}{cccc}
  1 & x_1 & x_2 & x_3 \\
  1 & a_1 & a_2 & a_3 \\
  1 & b_1 & b_2 & b_3 \\
  1 & c_1 & c_2 & c_3 \\
\end{array} \right| =0
    $$
    y hallar la ecuaci�n cartesiana del plano que pasa por los
    puntos $(1,2,1)$, $(-1,3,0)$, $(2,1,3)$.

    b) Escribir la ecuaci�n de una recta del plano que pasa por
    los puntos $(a,b)$, $(c,d)$ del plano y hallar la ecuaci�n
    cartesiana de la recta del plano que pasa por los puntos
    $(-1,-2)$, $(2,2)$.
    
    
    \n a) Siendo $a_0+a_1x+a_2x^2+a_3x^3=P(x)$ un polinomio de
    grado 3, hallar sus coeficientes para que
    $P(0)=2,\,P(1)=1,\,P(2)=-1,\,P(3)=0$.

Sol: $2+\frac{5}{6}x-\frac{5}{2}x^2+\frac{2}{3}x^3$.

    b)Demostrar que siempre se puede encontrar un polinomio de
    grado 3 que cumpla las condiciones $\{P(x_i)=y_i\}_{i=1}^4$
    siempre que todos los $x_1,\ldots,x_4$ sean distintos.

    c) Generalizar el resultado b): dados n+1 pares de puntos
    $\{x_i,y_i\}_{i=1}^{n+1}$ con $x_1,\ldots,x_{n+1}$ distintos,
    hay un �nico polinomio $P(x)$ de grado $n$ que cumple
    $P(x_i)=y_i,\,i=1,\ldots,n+1$.

    Obs�rvese que la parte c) nos indica c�mo hallar una funci�n
    polin�mica de grado $n$ que pase por $n+1$ puntos del plano
    siempre y cuando no haya dos puntos en la misma vertical.
    
    
    \n a) Obtener, en t�rminos de los determinantes de los
    menores diagonales de la matriz $A$,
    $$
    |A-\lambda I|=\left|\begin{array}{ccc}
      a_{11}-\lambda & a_{12} & a_{13} \\
      a_{21} & a_{22}-\lambda & a_{23} \\
      a_{31} & a_{32} & a_{33}-\lambda \\
    \end{array}
    \right|$$
    (Indicaci�n: descomponer las filas en sumandos con y sin
    $\lambda$).

    b) Aplicar la f�rmula obtenida para hallar los siguientes
    determinantes:
    $$
    |A-\lambda I|=\left|\begin{array}{ccc}
      1-\lambda & 2 & 1 \\
      0 & 3-\lambda & 1 \\
      3 & 1 & 2-\lambda \\
    \end{array}
    \right|,\hspace{1cm}\left|\frac{1}{3}\left(%
\begin{array}{ccc}
  2 & 1 & -2 \\
  -1 & -2 & -2 \\
  -2 & 2 & -1 \\
\end{array}%
\right)-\lambda I
    \right|    $$
    
    sol: $-{\lambda}^{3}  + {6 \lambda^{2} } - {7 \lambda} + 2$, $	
-{\lambda}^{3}  - \frac{1}{3}{\lambda}^{2}  + \frac{1}{3}\lambda + 1$

    \n Encontrar los valores de $a$ para que las siguientes
    matrices sean invertibles:
    $$
    a)\left(\begin{array}{ccc}
      a & 1 & 1 \\
      1 & a & 1 \\
      0 & 1 & a \\
    \end{array}\right),\hspace{1cm}b)\left(%
\begin{array}{ccc}
  -1 & -1 & 1 \\
  1 & 4 & -a \\
  -1 & a & 0 \\
\end{array}%
\right),\hspace{1cm}c)\left(%
\begin{array}{ccc}
  -1 & a & a \\
  a & -1 & a \\
  a & a & -1 \\
\end{array}%
\right)
    $$
Sol: a) $a\not=1$, $a\not=\frac{-1+\sqrt{5}}{2}$ y $a\not=\frac{-1-\sqrt{5}}{2}$  \quad b) $a\not=2$ y $a\not=-2$ \quad c) $a\not=-1$ y $a\not=\frac{1}{2}$.
    
    \n Sean matrices cuadradas $A\in M_{m\times m}$, $B\in M_{n\times n}$, $C\in M_{m\times
    n}$ y $D\in M_{(m+n)\times(m+n)}$ una matriz que se descompone
    por bloques como
    $$
    D=\left(%
\begin{array}{cc}
  A & C \\
  0 & B \\
\end{array}%
\right).
    $$

    a) Demostrar que $|D|=|A||B|$ utilizando la descomposici�n
    $$
    \left(%
\begin{array}{cc}
  A & C \\
  0 & B \\
\end{array}%
\right)=\left(%
\begin{array}{cc}
  I_m & 0 \\
  0 & B \\
\end{array}%
\right)\left(%
\begin{array}{cc}
  A & C \\
  0 & I_n \\
\end{array}%
\right)
    $$
    b) Demostrar, utilizando una descomposici�n similar, la misma
    igualdad $|D|=|A||B|$, en el caso
    $$
    D=\left(%
\begin{array}{cc}
  A & 0 \\
  C & B \\
\end{array}%
\right)
    $$
    (N�tese que ahora $C\in M_{n\times m}$).
    
    c) Calcula el determinante
$$\left| \begin{array}{cccccccc}
1 & -1 & 2 \\
   &  1  & 7 \\
   &      &  3 \\
   &      &      & 2 \\
   &      &      & 4  & -2 \\
   &      &      & 1  & 4   &  5   \\
   &      &      &     &      &      & 2 & 3 \\
   &      &      &     &      &      &    & 2  
\end{array}\right|$$

Sol: -240.

\n Estudiar si son independientes en su espacio vectorial correspondiente los siguientes conjuntos. En el caso en que sean linealmente dependientes, expresar uno de ellos como combinaci�n lineal de los restantes:

\apartado $\{(1, 2, 0), (0, 1, 1)\} \subset \mathbb{R}^3$, 

\apartado $\{(1, 2, 0), (0, 1, 1), (1, 1, 0)\} \subset \mathbb{R}^3$,

\apartado $\{(1, 2, 0), (0, 1, 1), (1, 1, -1)\} \subset \mathbb{R}^3$,

\apartado $\{(1, 2, 0), (0, 1, 1), (1, 1, 0), (3, 2, 1)\} \subset \mathbb{R}^3$,

\apartado $\{(i, 1 - i), (2, -2i - 2)\} \subset \mathbb{C}^2$,

\apartado $\{(i, 1, -i), (1, i, -1)\} \subset \mathbb{C}^3$.

\apartado $$\left\{\left(\begin{matrix} 0 & 1 \\ 0 & -2\end{matrix}\right),
\left(\begin{matrix}
1 & -1 \\ 
0 & 1 
\end{matrix}\right),
\left(\begin{matrix}
-2 & 5 \\ 
0 & -8
\end{matrix} \right) \right\} \subset \mathcal{M}_{2\times 2}(\mathbb{R})$$ 


\apartado Los polinomios $\{1, x - 2, (x - 2)^2, (x - 2)^3\} \subset P_{\mathbb{R}}^3[x]$.


\n \apartado Demostrar que si $\{f_1(x), f_2(x), \cdots , f_k (x)\}$ son funciones tales que existen $\{a_1 , a_2 ,\cdots, a_k \}$ verificando 
$$\left| \begin{array}{cccc}
f_1 (a_1 ) &f_2 (a_1 ) &  \cdots & f_k (a_1 ) \\
f_1 (a_2 )& f_2 (a_2 ) & \cdots & f_k (a_2 ) \\
\vdots & \vdots & & \vdots \\
f_1(a_k) & f_2(a_k) & \cdots & f_k(a_k) 
\end{array} \right| \not=0$$
son funciones linealmente independientes. 

\vspace{0.2cm}

Demuestra los siguientes enunciados utilizando el apartado anterior:

\apartado Los polinomios 
$\{1, x - 2, (x - 2)^2,\cdots, (x - 2)^n \} \subset P_{\mathbb{R}}^n[x]$ 
son linealmente independientes.

\apartado Para cualquier $a\in \mathbb{R}$, los polinomios  
$\{1, x - a, (x - a)^2, \cdots, (x - a)^n\} \subset P_{\mathbb{R}}^n[x]$ son linealmente independientes para cualquier $n\in \mathbb{N}$.


%\n Comprobar que los siguientes conjuntos de vectores son bases de los correspondientes espacios vectoriales: (son las llamadas bases can\'onicas) 

%\apartado $\{(1, 0, 0, \ldots , 0), (0, 1, 0, \ldots , 0), (0, 0, 1, \ldots , 0), \ldots , (0, 0, 0, \ldots , 1)\} \subset \mathbb{R}^n$. 

%\apartado $\left\{\left(\begin{matrix}
%1 & 0\\ 
%0  & 0\end{matrix} \right),
%\left(\begin{matrix}
%0 & 1\\
%0 & 0\end{matrix} \right),
%\left(\begin{matrix}
%0 & 0\\
%1 & 0\end{matrix} \right),
%\left(\begin{matrix}
%0 & 0 \\
%0 & 1 \end{matrix} \right)\right\}
%\subset \mathcal{M}_{2\times2} (\mathbb{R})$

%\apartado $\{1, x, x^2,\cdots , x^n\} \subset P^n_\mathbb{R} (x)$.

\vspace{0.2cm}

\n Estudiar si los siguientes conjuntos son bases de los espacios vectoriales dados: 

\apartado $\{(1, 0, 0, 1), (1, 0, 1, 0), (0, 1, 1, 0), (0, 1, 0, 1)\} \subset \mathbb{R}^4$

\apartado $\{(1, 0, 0, 1), (1, 0, 1, 0), (0, 2, 1, 0), (0, 1, 0, 1)\} \subset \mathbb{R}^4$.

\apartado $\left\{ \left(
\begin{matrix}
1 & 0\\
0 & 1\end{matrix} \right), 
\left( \begin{matrix}
1 &0 \\
1 &0\end{matrix} \right),
\left( \begin{matrix}
 0 &2 \\
1 &0\end{matrix} \right),
\left( \begin{matrix} 
0 &1 \\
0 & 1 \end{matrix}\right)
\right\}\subset \mathcal{M}_{2\times 2} (\mathbb{R})$.

\apartado $\left\{ \left(\begin{matrix}
1&  0\\ 
1 & 1
\end{matrix}\right),
\left( \begin{matrix}
0 &1\\ 
1 &1\end{matrix}\right),
\left( \begin{matrix}
1 &1\\ 
1 &2\end{matrix}\right),
\left(\begin{matrix}
0 &1\\
1 & 0\end{matrix}\right)\right.\subset \mathcal{M}_{2\times2} (\mathbb{R})$.

\vspace{0.2cm}

\n Encontrar bases de los siguientes subespacios vectoriales de 
$\mathcal{M}_{3\times3} (\mathbb{R})$:
 
\apartado Las matrices diagonales. 

\apartado Las matrices triangulares superiores. 

\apartado Las matrices triangulares inferiores. 

\apartado Las matrices sim\'etricas. 

\apartado Las matrices antisim\'etricas. 


\n \apartado Siendo $\{e_1 , e_2 , \ldots , e_n \}$ una base de $\mathbb{R}^n$, demostrar que los vectores: 
$$\begin{array}{cccccc}u_1 =  e_1,& 
u_2 =   e_1 + e_2 ,&
\cdots, &
u_{n-1} = e_1 + e_2 + \cdots + e_{n-1} ,&
u_n =  e_1 + e_2 + \cdots + e_{n-1}+e_n
\end{array}$$ 
son otra base de $\mathbb{R}^n$.

\apartado Hallar las coordenadas de los vectores $(n, n- 1, \ldots , 2, 1)$ y $(1, 2, \cdots , n - 1, n)$ de $\mathbb{R}^n$, en la base 
$\{u_1 , u_2 , \ldots, u_n \}$.

\vspace{0.2cm}

\n Hallar las coordenadas del vector $(3, 2, -1, 0)$ de $\mathbb{R} ^4$ 
en la base $$B=\{(1, 0, 0, 1), (1, 0, 1, 0), (0, 2, 1, 0), (0, 1, 0, 1)\}.$$
 
%\vspace{0.2cm}



\vspace{0.2cm}

    
   



\end{document}
